\section{The incidence identity and Hecke commutation}
\label{hec-section}

We prove the incidence identity for the actual affine matrix test
function and its passage to the theta columns. The scalar comes from
the change in the volume of the Hom lattice. The weights in the
global class sum are essential. All identities are first identities
of complex theta functions and their coefficient maps; the finite
coefficient-field construction then gives their arithmetic descent.

\subsection{Neighbors of a lattice with a residue functional}

Let $R$ be a discrete valuation ring with uniformizer $\varpi$,
fraction field $K_v$, and finite residue field $\mathbf k$ of order
$q$. A pair of rank $a$ consists of a free $R$-lattice $L\subset K_v^a$
and a nonzero $\mathbf k$-linear functional
$\ell:L/\varpi L\to\mathbf k$. We also write $\ell$ for its
composition with reduction on $L$. An upper neighbor of $(L,\ell)$ is
\[
 L^+=L+R\varpi^{-1}v,\qquad
 \ell^+(x)=\ell(\varpi x),\qquad v\in L,\quad\ell(v)=1.
 \tag{HC1}
\]
The displayed functional is well defined on $L^+/\varpi L^+$,
since $\varpi L^+\subset L$ and $\ell(\varpi y)=0$ for $y\in L$.
Its kernel in $L^+$ is exactly $L$.

\begin{lemma}\label{hec-neighbors}
There are $q^{a-1}$ distinct upper neighbors of $(L,\ell)$.
Every lower neighbor has underlying lattice $L^0=\ker\ell$.
Its functional $\ell^0:L^0/\varpi L^0\to\mathbf k$ is characterized
by
\[
             \ell^0(\varpi x)=\ell(x)\qquad(x\in L).
 \tag{HC2}
\]
There are $q^{a-1}$ such functionals, all nonzero.
\end{lemma}
\begin{proof}
Choose a basis $e_1,\ldots,e_a$ of $L$ with
$\ell(e_i)=0$ for $i<a$ and $\ell(e_a)=1$. The upper neighbors are
represented by $v=e_a+\sum_{i<a}u_i e_i$, $u_i\in\mathbf k$.
Indeed two such generators give the same lattice only if their
reductions span the same line in $L/\varpi L$; their last coordinates
are both one, so their reductions are equal. Formula (HC1) then fixes
the functional. The lattice $L^0$ has basis
$e_1,\ldots,e_{a-1},\varpi e_a$. Condition (HC2) prescribes the value
one on $\varpi e_a$ and leaves the other $a-1$ values arbitrary.
Conversely, for any such functional, the upper neighbor obtained
using $\varpi e_a$ is exactly $(L,\ell)$. The kernel identity in
(HC1) proves that every lower neighbor is of this form.
\end{proof}

For pairs $d=(L,\ell)$ and $d'=(M,\ell')$, of ranks $a,b$, and
$X\in\operatorname{Hom}_{K_v}(K_v^a,K_v^b)$, put
\[
 \phi_{d,d'}(X)=
  \mathbf1_{\{XL\subset M,\ \ell'X=\ell\text{ on }L\}}.
 \tag{HC3}
\]

\begin{proposition}\label{hec-local-count}
For every $X$,
\[
 \sum_{d_1:\ d'\to d_1}\phi_{d,d_1}(X)
       =\sum_{d_0:\ d_0\to d}\phi_{d_0,d'}(X).
 \tag{HC4}
\]
Both sums are either zero or one. Here arrows mean the upper-neighbor
relation (HC1).
\end{proposition}
\begin{proof}
Suppose $d'\to d_1=(M^+,\ell^+)$ and $\phi_{d,d_1}(X)=1$.
For $x\in L^0=\ker\ell$, the equality $\ell^+(Xx)=0$ gives
$Xx\in M$. Define $\ell^0(x)=\ell'(Xx)$ on $L^0$.
For $e\in L$ with $\ell(e)=1$,
\[
 \ell^0(\varpi e)=\ell'(\varpi Xe)=\ell^+(Xe)=1.
\]
More generally the same calculation gives (HC2), so $(L^0,\ell^0)$
is a lower neighbor and $\phi_{d_0,d'}(X)=1$.

Conversely suppose $d_0=(L^0,\ell^0)\to d$ and
$\phi_{d_0,d'}(X)=1$. Choose $e\in L$ with $\ell(e)=1$ and put
$v=X(\varpi e)\in M$. Then $\ell'(v)=1$. Define $M^+=M+R\varpi^{-1}v$
and its functional by (HC1). Since $L=L^0+Re$, the inclusion
$XL\subset M^+$ follows. For $x\in L^0$, $\ell^+(Xx)=0=\ell(x)$,
and for $x=e$ both values are one. Thus $\phi_{d,d_1}(X)=1$.

For uniqueness on the upper side, any solution satisfies
$Xe\in M^+$ and $\ell^+(Xe)=1$, hence $M^+=M+RXe$;
its functional is fixed by (HC1). On the lower side the underlying
lattice is $L^0$ and its functional must equal $\ell'X$.
The two constructions are inverse, which proves (HC4).
\end{proof}

For a self-dual Hermitian lattice over a CM field, changing the
component at a split prime $P$ forces the dual change at $\bar P$.
All other components and auxiliary frames remain fixed. Local patching
gives a unique global lattice, as in Lemma~\ref{acl-patching}.
Thus the same neighbor counts hold for the global pairs used here.
The test function is the product of (HC3) over the selected primes.
At a move at $P$, its other factors are identical on both sides of
(HC4); multiplying by them proves the global pointwise test identity.

\subsection{A finite theta calculation with a common reference lattice}

We explain the normalization in the lattice change, rather than
assuming that a finite identity of test functions automatically gives
an identity of differently normalized theta columns.
Fix positive even source and target ranks $a,b$.
Let $\mathcal S$ be a finite set of split primes, disjoint from its
conjugate. All period lattices in this paragraph agree outside
$\mathcal S\cup\overline{\mathcal S}$ and have dual components at
conjugate primes. Write $V_{\mathcal S}$ for the product of their
selected local vector spaces, with any one fixed additive Haar measure.
If $\mathcal L$ is one of these self-dual period lattices, write
$L=\prod_{P\in\mathcal S}\mathcal L_P$ for its selected component.
Write $v_{\mathcal L}$ for its normalized theta vector after removal
of the determinant character in (TE15). The exact lattice moves in
Theorem~\ref{te-equivariance} apply to these vectors.

For a compactly supported, locally constant function $\Phi$ on
$V_{\mathcal S}$, let $\mathcal T_{\mathcal L}(\Phi)$ be the finite
theta sum formed from the actual indicators (CS12a), with
$v_{\mathcal L}$ in place of $F$. Its support and its periods are
assumed to lie in the finite quotients used in that definition.
Increasing a support bound or refining a period does not change
this sum. The following proof includes that assertion.

\begin{lemma}\label{hec-common-theta}
For two such Hom period lattices $\mathcal L_1,\mathcal L_2$ and
one test function allowed for both,
\[
 \mathcal T_{\mathcal L_1}(\Phi)
   =\left(\frac{\operatorname{vol}L_2}
                  {\operatorname{vol}L_1}\right)^{1/2}
                                      \mathcal T_{\mathcal L_2}(\Phi).
 \tag{HC5}
\]
The square root is positive. This is an identity of entire functions.
\end{lemma}
\begin{proof}
Choose a sufficiently small selected Hom lattice $C$ contained in
$L_1\cap L_2$ and in the chosen periods of $\Phi$ for both lattices.
Choose a sufficiently
large selected Hom lattice $A$ containing $L_1,L_2$ and the support
of $\Phi$. These choices can be made within the Hom-lattice family:
enlarge the selected source and shrink the selected target to obtain
$C$, and do the reverse to obtain $A$. At conjugate primes impose
the dual changes and patch globally. Denote the resulting self-dual
period lattices by $\mathcal C,\mathcal A$.

Put $n_i=[L_i:C]$. The conjugate quotient dual to $L_i/C$ has order
$n_i$ and pairs perfectly with it by the Heisenberg commutator.
Consequently the projection of $v_{\mathcal L_i}$ onto the zero
selected coset modulo $C$ is
\[
 P_Cv_{\mathcal L_i}
       =\frac1{n_i}\sum_{y\in(L_i/C)^\vee}W_yv_{\mathcal L_i}
       =\frac1{\sqrt{n_i}}v_{\mathcal C}.
 \tag{HC6}
\]
The first equality is finite character orthogonality. The second is
the exact normalized lattice move, whose coefficient is
$n_i^{-1/2}$; moving the source and target consecutively gives the
same vector by the exact path independence proved in
Section~\ref{te-section}. This uses the positive normalization there.

Choose representatives of $A/C$ in $\mathcal A$ which are admissible
lifts for both original indicator systems. To do so, first choose
for each $i$ a positive denominator $c'_i$ clearing the selected
denominators of $A$ relative to $L_i$. At the conjugate selected
primes impose the additional divisibility by $c'_i$ required for
the half lattice in (CS12a), as well as membership in $\mathcal A$.
These are conditions at primes disjoint from the selected coset
conditions. The Chinese remainder theorem on $\mathcal A$ therefore
supplies all the required representatives simultaneously. For such a
representative $x$, the indicator of its coset applied to
$v_{\mathcal L_i}$ is
$n_i^{-1/2}W_{-x}v_{\mathcal C}$. These expressions are independent
of representatives: their difference lies in
$\mathcal A\cap\mathcal C$, the Weyl operators add on the integral
lattice $\mathcal A$, and this difference fixes $v_{\mathcal C}$.

For completeness, if $C\subset D\subset L_i$ is an intermediate
period (as for the actual scalar periods $cL_i$), the
identity of functions
$\mathbf1_{x+D}=\sum_{z\in D/C}\mathbf1_{x+z+C}$ gives the same
identity of these operators. Indeed the coefficient of a character
$y$ in the right side is multiplied by
\[
 \sum_{z\in D/C}e^{-2\pi iE(z,y)}
   =\begin{cases}|D/C|,&y\text{ is trivial on }D/C,\\
                 0,&\text{otherwise}.
     \end{cases}
\]
Together with the denominator $[L_i:C]$, this is exactly the
denominator and character sum for the coarser period $D$.
The lift independence in (CS12a) allows these common representatives
even when a support denominator exceeds the original period
denominator. Enlarging $A$ only inserts zero coefficients.
Thus the actual finite sum is
\[
 \mathcal T_{\mathcal L_i}(\Phi)
    =\frac1{\sqrt{n_i}}\sum_{x\in A/C}
                               \Phi(x)W_{-x}v_{\mathcal C}.
 \tag{HC7}
\]
Since $\operatorname{vol}L_i=n_i\operatorname{vol}C$, taking the
ratio proves (HC5). All sums are finite sums of translates of entire
theta functions, so the assertion holds as an identity of entire
functions, with no analytic convergence issue.
\end{proof}

Consider now ranks $a,b$ and lattices $N_1,N_2$ on the source and
$M_1,M_2$ on the target. They agree away from primes over a selected
rational prime $p$ which splits completely in $E$.
Thus $NP=p$ for each selected $P\mid p$, and $p$ is a local
uniformizer there. These are the actual selected-prime hypotheses.
For each such $P$, put
\[
 e_P=\operatorname{ord}_P\det N_2-\operatorname{ord}_P\det N_1,
 \qquad
 e'_P=\operatorname{ord}_P\det M_2-\operatorname{ord}_P\det M_1.
 \tag{HC8}
\]
The local period ideal is a unit. The determinant of
$\operatorname{Hom}(N,M)=M\otimes N^*$ therefore gives
\[
 \frac{\operatorname{vol}\operatorname{Hom}(N_2,M_2)_P}
      {\operatorname{vol}\operatorname{Hom}(N_1,M_1)_P}
                       =p^{b e_P-a e'_P}.
 \tag{HC9}
\]
This follows also from bases: an index $p^{e_P}$ change of $N$
changes the determinant of its dual $b$ times, and an index
$p^{e'_P}$ change of $M$ changes its determinant $a$ times.

The remaining character factor of the arithmetic column is
$D(N)/D(M)$. Away from the auxiliary level, its factor over a
prime of $F$ below $P$ is $w_P^{\operatorname{ord}_P\det N}$, where
$w_P$ is the nonzero unramified character value. If the selected
lift of that prime of $F$ is $\bar P$, both the determinant exponent
and the character are inverted, so the same factor results.
Here $w_{\bar P}=w_P^{-1}$ because the character is unramified at
both primes and trivial on the ground-field idele. Auxiliary frame
factors are unchanged under the present lattice move. Consequently
\[
 \frac{D(N_1)/D(M_1)}{D(N_2)/D(M_2)}
                         =\prod_{P\mid p}w_P^{e'_P-e_P}.
 \tag{HC10}
\]
Only the selected primes occur in this product.

\begin{proposition}\label{hec-kernel-move}
Let $\mathcal K_{N,M}(\Phi)$ be the actual theta coefficient map,
including the factors just described. Suppose the same test function
has the required support and periodicity for both pairs of lattices.
Then, over $\mathbb C$,
\[
 \mathcal K_{N_1,M_1}(\Phi)
   =\prod_{P\mid p}\left(w_P^{e'_P-e_P}
                    p^{(b e_P-a e'_P)/2}\right)
                                  \mathcal K_{N_2,M_2}(\Phi).
 \tag{HC11}
\]
Here $a,b$ are even, so the exponent is an integer.
\end{proposition}
\begin{proof}
Apply (HC5) and (HC9) to the determinant-normalized theta functions,
and then multiply by (HC10). The mixed derivatives defining the
column are linear operations on the same entire-function identity.
The polynomial coefficient extraction and the coefficient-module
maps are linear as well. Thus they preserve the scalar in (HC11).
No valuation estimate is substituted for this equality.
\end{proof}

\subsection{The pointwise column identity}

Fix pairs $d=(N,\ell)$ and $d'=(M,\ell')$, of even ranks $a,b$,
and a selected prime $P$. Write $N^0$ for the self-dual lattice
whose component at $P$ is $\ker\ell_P$, whose component at $\bar P$
is the dual one, and whose other components agree with $N$.
Let $\Phi_0$ be the prescribed auxiliary test function away from $p$.
The point kernel is
\[
          k(d,d')=\mathcal K_{N,M}(\phi_{d,d'}\Phi_0).
 \tag{HC12}
\]
All residue conditions at the other selected primes are included
in $\phi$. The auxiliary frames are fixed throughout each move.

\begin{theorem}\label{hec-point-identity}
Put $s=(b-a)/2$. Then
\[
 \sum_{d_1:\ d'\to d_1}k(d,d_1)
       =p^s\sum_{d_0:\ d_0\to d}k(d_0,d').
 \tag{HC13}
\]
\end{theorem}
\begin{proof}
If $d'\to d_1$ has underlying lattice $M_1$, the lower lattice
of $d_1$ is $M$. Change both lattices from $(N,M_1)$ to $(N^0,M)$.
At $P$ the two index exponents in (HC8) are $e_P=e'_P=1$,
and all other selected exponents vanish. Formula (HC11) therefore
has scalar exactly $p^{(b-a)/2}$; the character cancels.

The common test function is allowed for both pairs. If
$\phi_{d,d_1}(X)=1$, then $XN^0_P\subset M_P$ by the kernel
identity in (HC1). For periodicity, a period is divisible by $p$.
If $Z:N^0_P\to M_P$ is integral, then for $x\in N_P$ one has
$px\in N^0_P$ and
\[
                   pZx=Z(px)\in M_P=\ker\ell_{1,P}.
\]
Adding this period changes neither lattice inclusion nor the
residue equation. With the original pair, the same assertion
follows from $pM_{1,P}\subset\ker\ell_{1,P}$. Conditions at the
other selected primes are unchanged, and the lattices and frames
at the auxiliary primes are identical. Thus all support and
periodicity requirements in (HC11) hold. We obtain
\[
 k(d,d_1)=p^s\mathcal K_{N^0,M}(\phi_{d,d_1}\Phi_0).
\]
Every lower neighbor $d_0\to d$ has underlying lattice $N^0$,
so $k(d_0,d')=\mathcal K_{N^0,M}(\phi_{d_0,d'}\Phi_0)$.
Both sums are finite. Use linearity and (HC4) to conclude.
\end{proof}

\subsection{The weights in the global class sum}

Let a group $\Gamma$ act on a set $\mathscr X$ with finitely many
orbits and finite stabilizers. Let $\mathscr R\subset\mathscr X^2$
be an invariant relation with finite fibers in both directions.
Choose one representative $x_c$ of each orbit and put
$w_c=|\Gamma_{x_c}|^{-1}$.

\begin{lemma}\label{hec-incidence-mass}
For a simultaneous $\Gamma$-invariant vector-valued function $F$
on $\mathscr R$,
\[
 \sum_c w_c\sum_{y:\,(x_c,y)\in\mathscr R}F(x_c,y)
       =\sum_c w_c\sum_{x:\,(x,x_c)\in\mathscr R}F(x,x_c).
 \tag{HC14}
\]
\end{lemma}
\begin{proof}
There are finitely many $\Gamma$-orbits of related pairs, since
there are finitely many first-coordinate orbits and finite fibers.
For one such orbit, choose $(x,y)$ with first coordinate one of
the representatives. Its contribution to the left side is
\[
 \frac1{|\Gamma_x|}[\Gamma_x:\Gamma_x\cap\Gamma_y]F(x,y)
                 =\frac{F(x,y)}{|\Gamma_x\cap\Gamma_y|}.
\]
Computing the same orbit using the second coordinate gives exactly
the same value. Summing over these finitely many orbits proves (HC14).
\end{proof}

The coefficient spaces over $\mathscr X$ may carry representations
of $\Gamma$. Equivariant sections and equivariant kernels are allowed:
the expression $k(x,z)f(y)$ is simultaneously invariant in $(x,y)$,
because the source transformation of the kernel cancels that of $f$.
This is the only invariance required in (HC14).

\begin{theorem}\label{hec-global-identity}
On the actual finite spaces of equivariant forms, let $U_a,U_b$
sum over upper neighbors, and define the theta column by
\[
              (Kf)(d')=\sum_c w_c\,k(d_c,d')f(d_c).
 \tag{HC15}
\]
Then
\[
                  U_bK=p^{(b-a)/2}KU_a,
 \qquad U_b^rK=p^{r(b-a)/2}KU_a^r\quad(r\ge0).
 \tag{HC16}
\]
\end{theorem}
\begin{proof}
The class interpretation and its stabilizer weights were proved in
Sections~\ref{acl-section} and \ref{acc-section}. The local neighbors
there identify with (HC1); the residue functional records the
quotient map, so no further multiplicity occurs. Finiteness of the
fibers follows from Lemma~\ref{hec-neighbors}. Exact unitary
covariance in (TE17), together with the natural changes of auxiliary
frames, gives the equivariance of this kernel and of the forms.

Expand the left side of (HC16) using (HC15), interchange the two
finite sums, and use (HC13). The result, after division by $p^s$, is
\[
 \sum_c w_c\sum_{d_0:\,d_0\to d_c} k(d_0,d')f(d_c).
\]
Apply (HC14) with $F(x,y)=k(x,d')f(y)$. It equals
$\sum_c w_c k(d_c,d')\sum_{y:\,d_c\to y}f(y)=(KU_af)(d')$.
This proves the first equality, and induction proves the second,
including $r=0$.
\end{proof}

Proposition~\ref{jcf-common-field} constructs the number field $K$
containing the needed mixed derivatives, character values, and frame
coefficients for this fixed finite construction. All the coefficient
maps are therefore defined over $K$, and (HC16) descends to $K$ by
injectivity of $K\hookrightarrow\mathbb C$. This uses the finite
field construction, not merely the algebraicity of individual jets.
It discharges the commutation hypothesis in the finite valuation
iteration of Section~\ref{cv-section}. The exact source definitions
and the finite checks are identified in the accompanying trace
and manifest \cite{alpo-uniform-source}.

\begin{example}[Why the stabilizer weights are required]
Let $S_3$ act on one fixed point $*$ and on three permuted points.
Relate $*$ to each of the other three points and set $F=1$ on these
three pairs. There are two orbits of points. The two unweighted
class sums in (HC14) would be $3$ and $1$. The actual weights give
$3/6=1/2$ and $1/2$. Thus dropping the stabilizers would invalidate
the passage to class sums. They are present in the submitted proof.
\end{example}


